% =====================================================================
% Chapter 1 -- Introduction
%
% Redrafted from 2026-10-23 (outline: outlines/ch1_introduction.md, revision 2).
% The author's old notes are kept, commented out, at the end of this file
% until all sections are redrafted (backup: backups/introduction_2026-10-23_before_redraft.tex).
% =====================================================================

\chapter{Introduction}
\label{chap:intro}

% ---------------------------------------------------------------------
\section{Motivation}
\label{sec:intro-motivation}

% CLAUDE-DRAFT
Medical image segmentation assigns every voxel of an image either to a structure, such as an organ or a tissue, or to the background.
The voxels assigned to a structure form its mask.
Masks can be drawn by hand, but this is laborious, because in a three-dimensional (3D) volume the structure has to be outlined on every slice in which it appears.
Different readers also draw them differently~\cite{desai2021skmtea}.

% CLAUDE-DRAFT
Automatic segmentation with deep learning has commonly meant training one model for each dataset and its target structures, for example with nnU-Net~\cite{isensee2018nnunetselfadaptingframeworkunetbased}.
Universal segmentation models replace these per-dataset models with one network that segments many structures.
For example, the developers of the universal model SAT compared it with 72 nnU-Nets, one trained on each of their datasets~\cite{zhao2025sat}.
The structure to segment is specified by an additional input, the prompt~\cite{zhao2025sat,rokuss2025voxtellfreetextpromptableuniversal}.
The prompt can be spatial, such as a point or a box marked in the image~\cite{Ma_2024,cheng2023sammed2d}, or text that names the target.

% CLAUDE-DRAFT
Text-promptable models differ in the text they accept.
SAT is prompted with terms from a fixed vocabulary~\cite{zhao2025sat}.
In contrast, VoxTell accepts both a single anatomical term and a full clinical sentence~\cite{rokuss2025voxtellfreetextpromptableuniversal}.
We call models of this kind \emph{free-text promptable}.
In both models, the structure to segment is not tied to an output channel of the network.
The same weights serve every structure, and the text prompt alone selects the target~\cite{zhao2025sat,rokuss2025voxtellfreetextpromptableuniversal}.

% CLAUDE-DRAFT
This thesis adapts VoxTell, which was trained on more than 62{,}000 \gls{ct}, \gls{mri} and \gls{pet} volumes~\cite{rokuss2025voxtellfreetextpromptableuniversal}.
On datasets it was not trained on, it reached a mean \gls{dsc} of 70.85\,\%, where the \gls{dsc} measures the overlap between the predicted and the reference mask.
On knee \gls{mri} from the SKM-TEA dataset~\cite{desai2021skmtea}, its mean \gls{dsc} was 10.72\,\%, averaged over four cartilages and two menisci~\cite[Suppl.~Table~9]{rokuss2025voxtellfreetextpromptableuniversal}.
The VoxTell authors attribute this low value to the small share of knee images in the training data and to structures missing from it~\cite[Sec.~6]{rokuss2025voxtellfreetextpromptableuniversal}.
As future work, they suggest adapting the model with a few labelled examples.

% CLAUDE-DRAFT
Automatic segmentation of the knee cartilages and menisci is useful because quantitative biomarkers for osteoarthritis are measured within these structures~\cite{desai2021skmtea}.
One example is the T2 relaxation time, which reflects the degeneration of collagen-rich tissue.
These values can only be read out once the structures are outlined.
The SKM-TEA authors name manual outlining as one obstacle to the use of quantitative \gls{mri}.

% CLAUDE-DRAFT
With few manual outlines available, VoxTell would have to learn the knee structures from a few labelled volumes.
New structures may also become relevant over time, and the model would then have to adapt to each of them in turn.
Retraining VoxTell from scratch on its training data together with each new structure would be expected to preserve most of what it could segment before.
However, training the published model took about six days on 64 NVIDIA A100 GPUs~\cite{rokuss2025voxtellfreetextpromptableuniversal}.

% ---------------------------------------------------------------------
\section{Problem Statement}
\label{sec:intro-problem}

% CLAUDE-DRAFT
Instead of being retrained from scratch, the model can be adapted to one new structure at a time.
We call the learning of one new structure, named by one text prompt, a \emph{task}.
Continual learning is the setting in which a model learns from a sequence of tasks, using the data of the current task and limited or no data of earlier tasks~\cite[Sec.~2]{delange2022}.
After each task, the model has to segment the new structure while keeping the structures learned earlier in the sequence.
It also has to keep what it could segment before the first task, which we call its pretrained ability (\cref{fig:ch1_setting}).

% CLAUDE-DRAFT
The simplest way to learn a task is to fine-tune the model on its data.
However, training on later tasks can lower the performance on earlier ones, an effect known as catastrophic forgetting~\cite[Sec.~1]{delange2022}.
For pretrained vision--language models, continual fine-tuning can also lower the zero-shot accuracy on data the model was not trained on~\cite{zheng2023preventingzeroshottransferdegradation}.
In our setting, this corresponds to a loss of pretrained ability.
The more the model is allowed to change, the better it can learn the new structure, but the more it may lose of what it knew.
De Lange et al.\ call this conflict the stability--plasticity dilemma, with stability meaning the retention of previous knowledge and plasticity the integration of new knowledge~\cite[Sec.~1]{delange2022}.

\begin{figure}[tb]
\centering
\includegraphics{figures/ch1_setting.pdf}
% CLAUDE-DRAFT
\caption[Setting of this thesis]{Setting of this thesis.
The pretrained VoxTell segments the lung well but the knee structures poorly (\gls{dsc} from our evaluation; for the knee, the mean over its six structures).
It is then adapted to the six knee structures one after another, each named by a text prompt and learned from a few labelled volumes.
Tasks~3 to~5 are not shown.
After every task, the same model has to segment all of these structures, with the lung standing for its pretrained ability.
Images: example volumes of \cref{fig:ch4_examples} with their reference masks.}
\label{fig:ch1_setting}
\end{figure}

% CLAUDE-DRAFT
With few labelled volumes per task, it is not clear in advance how the balance between stability and plasticity shifts.
Fewer volumes mean fewer updates when the number of training epochs is fixed.
Fewer updates give the model less opportunity to learn the new structure, but they may also change less of what it knew.
Here, \emph{few-shot} means training the model on $k$ labelled volumes per task, one shot being one labelled training volume.

% CLAUDE-DRAFT
Existing continual-learning results do not show how few labelled volumes shift the balance between stability and plasticity, because they were obtained in other settings.
The methods for pretrained vision--language models reviewed in \cref{sec:cl-vlm} were developed and evaluated on classification or on tasks with text output.
In medical image segmentation, many continual-learning methods give each new task or class components of its own, such as a decoder or an output head~\cite{ji2023continualsegmentsingleunified,guo2025continuallearningdrivenmodelaccurate,zhang2023continuallearningabdominalmultiorgan,ye2025medseqft}.
VoxTell needs no such components, because a new structure is named by a new text prompt and segmented with the same weights.

% CLAUDE-DRAFT
In our literature search, we found no work that adapts a pretrained free-text promptable 3D segmentation model to a sequence of new structures with few labelled volumes per task (\cref{sec:sota-summary}).
The closest works use a two-dimensional (2D) model with full training sets~\cite{gao2026medcrpcl,gao2025paewc,gao2026medpeftcl} or a 3D model without text prompts~\cite{chen2026fewshotcontinuallearning3d}, and none of them measures how much of the pretrained ability is kept.
Other works measure it for CLIP~\cite{Zhang_2024_CVPR} and for segmentation models prompted with points or boxes~\cite{wang2025casam,wang2024sammed3dmoe}, and our control task builds on this idea.

% ---------------------------------------------------------------------
\section{Research Question and Objectives}
\label{sec:intro-rq}

% CLAUDE-DRAFT
This thesis addresses the following research question:
\begin{quote}
\textbf{How do different continual-learning strategies trade off adaptation performance and catastrophic forgetting when adapting free-text promptable universal 3D medical image segmentation models under few-shot conditions?}
\end{quote}
By \emph{adaptation performance} we mean how well the model segments a new structure right after it has been trained on it.
In the research question, \emph{forgetting} covers the loss on earlier tasks of the sequence as well as the loss of pretrained ability.
Few-shot conditions are $k = 5$, 10 and 20 labelled volumes per task, which we compare with the full training split of 108 volumes per task.

% CLAUDE-DRAFT
We pursue four objectives:
\begin{enumerate}
  \item Develop an evaluation protocol for the setting of the research question that measures adaptation, the forgetting of earlier tasks and the loss of pretrained ability separately.
  \item Compare continual-learning strategies with each other and with fine-tuning on each task in turn, and determine how they trade adaptation against forgetting.
  \item Investigate how the number of labelled volumes per task influences adaptation and forgetting, separating the effect of fewer volumes from that of fewer updates.
  \item Analyse how the order of the tasks and the design choices within the strategies influence adaptation and forgetting.
\end{enumerate}

% ---------------------------------------------------------------------
\section{Approach and Scope}
\label{sec:intro-approach}

% CLAUDE-DRAFT
We study the research question with VoxTell~\cite{rokuss2025voxtellfreetextpromptableuniversal} as the pretrained model and six structures of the SKM-TEA knee \gls{mri} dataset~\cite{desai2021skmtea} as tasks.
These are the patellar, femoral, medial tibial and lateral tibial cartilage and the medial and lateral meniscus.
The tasks are trained in two task orders, an original order chosen at random and its reverse, to see whether the results depend on the order.
The lung in the chest \gls{ct} volumes of AeroPath~\cite{støverud2023aeropathairwaysegmentationbenchmark} serves as a control task.
It is not trained on during the adaptation and is used only to measure how much of the pretrained ability is kept.
After each task, the model is tested on all six tasks and on the control task (\cref{sec:method-setting-protocol}).

% CLAUDE-DRAFT
The baseline is sequential fine-tuning, which trains all parameters on each task and adds nothing to reduce forgetting.
Three strategies are based on \gls{lora}~\cite{hu2022lora}.
\gls{lora} keeps the pretrained weights frozen and trains only a low-rank update of them, called the adapter.
The number of parameters of the adapter grows with its rank~$r$.
Task-specific LoRA trains one adapter of rank $r = 4$ per task and selects it by the task identity at inference.
Shared LoRA trains one adapter of rank $r = 1024$ for all tasks, and shared LoRA with replay adds a memory of training samples from earlier tasks.
The prompt-tuning strategy (adapted from CPE-CLIP~\cite{D_Alessandro_2023}) freezes the pretrained weights and trains learned prompt tokens.
It also trains a meta-net, a small network that shifts the encoded text prompt.
The meta-net is our addition, inspired by CoCoOp~\cite{zhou2022cocoop}.
The distillation-based strategy (adapted from ZSCL~\cite{zheng2023preventingzeroshottransferdegradation}) trains all parameters and distils the features of the frozen pretrained model into the adapted one.
Together, the strategies cover replay, regularisation and parameter-efficient adaptation, the three families of continual-learning methods for vision--language models~\cite[Sec.~4]{liu2026continuallearningvlmssurvey}.
All strategies except task-specific LoRA use one parameter set for all tasks.

% CLAUDE-DRAFT
All strategies are trained with the full training split.
Sequential fine-tuning, task-specific LoRA and shared LoRA are also trained with $k = 5$, 10 and 20 volumes per task.
With a fixed number of epochs, fewer volumes also mean fewer updates.
Further runs keep the number of updates of the full training split and so separate the two effects (\cref{subsec:exp-data}).\todo{defense: the research question names few-shot conditions, but the distillation-based and prompt-tuning strategies are run only with the full training split (D46)}

% CLAUDE-DRAFT
From the test results, we compute three primary metrics (\cref{sec:bg-metrics}).
Adaptation is the \gls{dsc} on each task right after training on it, averaged over the six tasks.
Final Forgetting is the mean drop of the \gls{dsc} on the earlier tasks from its best earlier value to its value after the last task~\cite{Chaudhry_2018}.
The control-task drop is the decrease of the \gls{dsc} on the control task from the pretrained model to the model after the last task.
Together, Final Forgetting and the control-task drop measure retention, the performance kept on earlier tasks and on the control task.
We also report two secondary metrics.
Final Performance is the mean \gls{dsc} on all tasks after the last task, and Transfer is the mean \gls{dsc} on each task before it is trained~\cite{zheng2023preventingzeroshottransferdegradation}.
The pretrained model is evaluated in the same way and serves as a reference.
As an upper reference, joint fine-tuning trains the model on all six tasks at once.
It was run with the full training split and with each~$k$.

% CLAUDE-DRAFT
The study uses one model, one dataset for the tasks, one control task, one fixed text prompt per structure and one draw of training volumes per~$k$.
Each configuration was trained once, so the only variability we can estimate is that over the test volumes (\cref{subsec:exp-metrics}).
The conclusions are therefore limited to this setting (\cref{sec:exp-limitations}).
Clinical deployment and studies with clinical users are outside the scope of this thesis, but we report the training cost of the strategies (\cref{sec:exp-compute}).

% ---------------------------------------------------------------------
\section{Contributions}
\label{sec:intro-contributions}
% Contributions decided by the author on 2026-10-26 (option A, thesis_notes.md §1); one per objective since 2026-11-01.

% CLAUDE-DRAFT
The contributions of this thesis are:
\begin{enumerate}
  \item An evaluation protocol that measures the loss of pretrained ability separately from the forgetting of earlier tasks, through a held-out control task, with the pretrained model and joint fine-tuning as reference points (\cref{sec:method-setting-protocol}).
  \item A comparison of sequential fine-tuning and five continual-learning strategies adapted to VoxTell under one shared training configuration, together with a diagnosis of why task-specific LoRA did not learn patellar cartilage (\cref{sec:method-strategies,sec:exp-main,sec:exp-patellar}).
  \item An analysis of how the amount of labelled data affects adaptation and forgetting, with three of the strategies and joint fine-tuning, which separates fewer volumes from fewer updates by training on five volumes with as many updates as on the full training split (\cref{sec:exp-few-shot}).
  \item An analysis of how the task order and the rank of the LoRA adapters affect adaptation and forgetting, with two task orders and ranks from $r = 4$ to $r = 2048$ (\cref{sec:exp-rank,sec:exp-task-order}).
\end{enumerate}

% CLAUDE-DRAFT
Sequential fine-tuning, LoRA, replay~\cite[Sec.~3.1]{delange2022}, ZSCL and CPE-CLIP are existing methods.
The work of this thesis is their application to VoxTell with the changes described in \cref{sec:method-strategies}, the evaluation protocol and all experiments.
The works reviewed in \cref{sec:sota-summary} do not study the setting of \cref{sec:intro-problem}.

% ---------------------------------------------------------------------
\section{Thesis Outline}
\label{sec:intro-outline}

% CLAUDE-DRAFT
\Cref{chap:sota} reviews text-promptable medical image segmentation and few-shot adaptation of segmentation models, as well as continual learning for pretrained vision--language models and for medical image segmentation.
It ends with the gap that this thesis addresses.
\Cref{chap:background} introduces VoxTell, states the continual-learning problem in formal terms, describes the published methods on which the strategies are based and defines the evaluation metrics.
\Cref{chap:data} describes the SKM-TEA and AeroPath data, the tasks and the data splits.
\Cref{chap:method} defines the evaluation protocol, the adaptation of VoxTell to one task and the strategies with their complexity.
\Cref{chap:experiments} gives the experimental setup and the results.
It compares the strategies, diagnoses why task-specific LoRA did not learn patellar cartilage, and reports the experiments on the amount of labelled data, the LoRA rank and the task order, as well as the training cost, and ends with the limitations of the experiments and the results that may generalise beyond them.
\Cref{chap:summary} summarises the results, answers the research question and outlines future work.
The appendix lists the training and test volumes and the parameters of all runs (\cref{sec:app-volumes,sec:app-params}), reports the additional evaluations, including two components adopted from ZSCL and CPE-CLIP (\cref{sec:app-additional,sec:exp-components}), and gives the evaluation matrices of all runs and the training loss of shared LoRA (\cref{sec:app-matrices,sec:app-shared-loss}).

% =====================================================================
% AUTHOR'S OLD NOTES (2026-10-22 and earlier), not part of the thesis.
% Delete once all sections above are drafted.
% =====================================================================
% \chapter{Introduction}
%
% Medical image segmentation aims to delineate anatomical structures or pathological regions and produce masks that support the analysis of anatomy, disease, and treatment-related findings. Deep learning-based segmentation has progressed from task-specific models developed for individual datasets towards text-promptable universal models designed to support segmentation across multiple structures, modalities, and tasks \cite{zhao2025sat, Zhao_2024, text3dsam, rokuss2025voxtellfreetextpromptableuniversal}. However, text-promptable universal models can still exhibit reduced performance on previously unseen tasks, particularly when the target task or imaging characteristics differ substantially from those represented during pretraining \cite{rokuss2025voxtellfreetextpromptableuniversal, zhao2025sat}. Adapting the model to such tasks therefore becomes necessary.
%
% This thesis investigates this problem from a continual-learning perspective, focusing on free-text promptable universal models for three-dimensional (3D) medical image segmentation. The central challenge is to adapt the model to newly introduced tasks while preserving its performance on previously learned tasks, particularly when only a small number of labelled volumes are available. The following sections establish the background and problem formulation for this study.
%
% \section{Universal Models for Medical Image Segmentation}
%
% Three-dimensional medical imaging is obtained using modalities such as computed tomography (CT), magnetic resonance imaging (MRI), and positron emission tomography (PET), which commonly produce volumetric data composed of spatially aligned slices. Unlike two-dimensional images, these volumes preserve spatial context across adjacent slices, which must be considered when segmenting anatomical structures or pathological regions \cite{DAYARATHNA2024103046}. 
%
% Deep learning is a widely used approach for automated medical image segmentation. Convolutional neural networks (CNNs), including U-Net-based architectures, have been extensively adapted to 3D medical imaging \cite{3d_unet, li2018hdenseunethybriddenselyconnected, wang2020nonlocalunetbiomedicalimage}. Many of these approaches are specialised for a particular combination of imaging modality, dataset, and target structure. The nnU-Net framework provides a systematic implementation of this task-specific paradigm by automatically adapting its configuration and training pipeline to a given dataset \cite{isensee2018nnunetselfadaptingframeworkunetbased}. Although task-specific models can achieve strong segmentation performance within their intended settings, deploying separate models for an increasing number of anatomical structures, modalities, and datasets becomes less practical \cite{zhao2025sat}.
%
% Universal medical image segmentation models aim to address this limitation by allowing a common model to produce different segmentation outputs for different target structures or regions specified through an input prompt. Early approaches to promptable universal segmentation models relied primarily on visual interaction mechanisms, requiring users to provide spatial prompts, such as points or bounding boxes, to specify the segmentation target \cite{Ma_2024, du2025segvoluniversalinteractivevolumetric, cheng2023sammed2d}. Text-promptable segmentation \cite{carion2026sam3segmentconcepts} offers clinicians a more intuitive alternative by allowing them to extract anatomical structures or pathological regions using simple natural language descriptions, such as 'segment the left ventricle' and obtain the corresponding mask. 
%
% However, medical imaging systems deployed in practice may encounter changes in data distributions as new imaging conditions, anatomical targets, pathologies, or acquisition settings are introduced \cite{qazi2026clmedicalimaging}. When multiple new tasks are introduced sequentially, adapting the model can also affect its performance on previously learned tasks \cite{kemker2017measuringcatastrophicforgettingneural}. This challenge motivates the continual-learning setting considered in this thesis.
%
%
% \section{Continual and Few-Shot Learning}
%
% One straightforward adaptation strategy is to fine-tune the pretrained model using data from the newly introduced task. However, sequential fine-tuning can cause catastrophic forgetting, where performance on previously learned tasks can decrease after subsequent updates \cite{french1999catastrophic}. For multimodal models, catastrophic forgetting can be particularly severe, with different modalities exhibiting different rates of forgetting and interactions between modalities further complicating continual adaptation \cite{yu2026recentadvancesmultimodalcontinual}.
%
% Continual learning (CL) addresses this problem by considering learning as a sequence of updates rather than a single training stage \cite{van_de_Ven_2025}. At each step, a model receives data associated with a newly introduced task and is expected to learn the corresponding capability while retaining knowledge acquired from earlier tasks. This requires balancing stability, the ability to preserve previously learned knowledge, and plasticity, the ability to acquire new knowledge. Several methods have been proposed to mitigate catastrophic forgetting in unimodal models, including regularisation-based \cite{Kirkpatrick_2017}, rehearsal-based \cite{icarl2017}, and dynamic architecture-based \cite{yoon2018lifelonglearningdynamicallyexpandable} approaches. However, applying continual learning to multimodal foundation models introduces additional challenges due to their large scale and architectural complexity \cite{liu2026continuallearningvlmssurvey}. Continual learning can be formulated under different settings, including task-incremental, class-incremental, and domain-incremental learning, which differ in how task identity, class information, and data distributions are treated across learning steps \cite{liu2026continuallearningvlmssurvey}.
%
% The continual adaptation problem becomes more challenging when only a limited number of labelled examples are available for each newly introduced task. In medical imaging, obtaining high-quality voxel-level annotations requires domain expertise and can involve substantial effort, particularly for three-dimensional volumes \cite{Litjens_2017}. Few-shot learning refers to the adaptation of a model using only a small number of labelled examples from the target task \cite{theekshana26fewshot}. Few-shot conditions are used to evaluate how continual-learning strategies behave under limited supervision.
%
% \section{Problem Statement and Motivation}
% Early work on continual learning in medical image segmentation has largely considered task-specific or relatively compact models \cite{karani2018lifelonglearningapproachbrain, memmel2021adversarialcontinuallearningmultidomain, ranem2022continualhippocampussegmentationtransformers, gonzalez2023lifelong}. Modern free-text promptable universal segmentation models, in contrast, combine large-scale visual and language representations with cross-modal interaction components. Their architectural complexity and multimodal nature introduce additional considerations for continual adaptation \cite{liu2026continuallearningvlmssurvey} that are not necessarily captured by approaches developed for conventional segmentation models.
%
% Although recent studies have begun to investigate continual learning for medical foundation models \cite{chen2026fewshotcontinuallearning3d, Zhang_2025, chen2025forgingdynamicmemoryretrievalguided}, a systematic comparison of continual-learning strategies for free-text promptable universal 3D medical image segmentation remains limited. In particular, the behaviour of different continual-learning strategies in this setting remains insufficiently understood. This question is especially relevant under few-shot conditions.
%
% \section{Research Question and Objectives}
%
% Based on the problem formulated in the previous section, this thesis investigates the following research question:
%
% \medskip
% \noindent\textbf{How do different continual-learning strategies trade off adaptation performance and catastrophic forgetting when adapting free-text promptable universal 3D medical image segmentation models under few-shot conditions?}
%
% \medskip
% \noindent The research question is addressed through the following objectives:
% \begin{enumerate}
%     \item \textbf{Develop a few-shot continual-learning evaluation protocol}
%     for free-text promptable universal 3D medical image segmentation.
%
%     \item \textbf{Compare different continual-learning strategies} under the
%     proposed evaluation protocol in terms of their adaptation and forgetting
%     behaviour.
%
%     \item \textbf{Investigate the influence of the amount of labelled training
%     data} on the performance of the evaluated strategies under few-shot
%     conditions.
%
%     \item \textbf{Analyse the influence of experimental factors} on adaptation and forgetting through quantitative experiments and ablation studies.
% \end{enumerate}
%
% \section{Contributions}
%
% The main contributions of this thesis are:
%
% \begin{enumerate}
%     \item \textbf{A systematic empirical evaluation of continual adaptation for free-text promptable universal 3D medical image segmentation.}
%
%     \item \textbf{A comparative analysis of adaptation, forgetting, transfer, and retention under limited supervision,} providing an empirical basis for understanding the behaviour of continual-learning approaches in this setting.
%
%     \item \textbf{An investigation of factors affecting continual adaptation performance,} including the amount of available labelled data, task ordering, and adaptation configuration.
%
% \end{enumerate}
%
% \section{Scope and Delimitations}
%
% The scope of this thesis is defined by the model, datasets, continual-learning
% protocol, adaptation strategies, and evaluation measures considered in the experiments. The study is limited to the VoxTell \cite{rokuss2025voxtellfreetextpromptableuniversal} model and the datasets \cite{desai2021skmtea, støverud2023aeropathairwaysegmentationbenchmark} included in the defined experimental protocol. 
%
% The continual-learning experiments are conducted using the task sequences and few-shot conditions specified in the experimental design. The study compares selected continual-learning strategies that employ different mechanisms for adaptation and mitigating catastrophic forgetting, including parameter-efficient adaptation, prompt-based adaptation, and regularisation. In addition, the study evaluates the combination of parameter-efficient adaptation with replay. The study does not aim to develop a new continual-learning algorithm or to establish a universally optimal strategy for medical image segmentation.
%
% Evaluation is based on segmentation performance and continual-learning measures that quantify adaptation to newly introduced tasks and forgetting of existing capabilities. Factors such as clinical deployment, inference latency, computational cost in clinical environments, and prospective validation with clinical users are outside the scope of this study. The experiments are conducted within the available computational and storage resources, which constrain the model configurations and experimental settings considered.
%
% Finally, the conclusions are limited to the model, datasets, task sequences, adaptation methods, and few-shot conditions evaluated in this work; their generalisability to other medical imaging tasks or foundation models remains to be established.
%
%
% % The remainder of this thesis is organised as follows. Chapter 2 reviews the state of the art in universal medical image segmentation, continual learning, and related adaptation approaches. Chapter 3 presents the theoretical background required to understand the methods and evaluation used in this thesis. Chapter 4 describes the datasets used in the experiments. Chapter 5 presents the experimental setup, results, and analysis of the investigated continual-adaptation strategies. Finally, Chapter 6 summarises the main findings of the thesis, discusses its limitations, and outlines directions for future work.
